\chapter{O Sacramento da Ordem}

O sacramento da Ordem é o sacramento pelo qual a missão confiada por Cristo a seus apóstolos continua a ser exercida na Igreja até o fim dos tempos: é, pois, o sacramento do ministério apostólico. Compreende três graus: o episcopado, o presbiterado e o diaconato. Não é um privilégio pessoal nem uma honraria institucional; é um serviço --- uma diaconia --- pelo qual Cristo, Cabeça do Corpo, age visivelmente no meio da sua Igreja, ensinando, santificando e governando o povo de Deus. Nos parágrafos 1536 a 1600, o Catecismo da Igreja Católica expõe a natureza, a história, os graus, a celebração e os efeitos deste sacramento, situando-o dentro da única sacralidade de Cristo, de quem tudo provém e a quem tudo retorna.

\section{Por que se chama sacramento da ``Ordem''? (1536--1538)}

O ponto de partida do Catecismo é a própria palavra ``Ordem''. Na Roma antiga, o termo \textit{ordo} designava um corpo civil constituído, especialmente um corpo governante. \textit{Ordinatio} significava a incorporação a um \textit{ordo}. Na Igreja, desde os tempos mais antigos, existem corpos constituídos que a Tradição --- com fundamento na Sagrada Escritura --- sempre chamou de \textit{taxeis} em grego ou \textit{ordines} em latim: daí a liturgia falar do \textit{ordo episcoporum}, do \textit{ordo presbyterorum}, do \textit{ordo diaconorum}. Outros grupos também recebem esse nome: catecúmenos, virgens, cônjuges, viúvos... A inserção em um desses corpos era realizada por um rito chamado \textit{ordinatio}, ato religioso e litúrgico que constituía uma consagração, uma bênção ou um sacramento.

Hoje, a palavra ``ordenação'' é reservada ao ato sacramental que integra um homem na ordem dos bispos, dos presbíteros ou dos diáconos, indo além de uma simples eleição, designação, delegação ou instituição pela comunidade: confere um dom do Espírito Santo que permite o exercício de um ``poder sagrado'' (\textit{sacra potestas}) que só pode vir do próprio Cristo, pela sua Igreja. Por isso, a ordenação também recebe o nome de \textit{consecratio}: é uma separação e um investidura pelo próprio Cristo para a sua Igreja. O sinal visível desta ordenação é a imposição das mãos pelo bispo, acompanhada da oração consecratória.

O sacramento da Ordem é, portanto --- já nesta definição inicial ---, inteiramente cristológico e eclesial: não é o ordinando que se apodera de um poder; é Cristo que chama, que separa, que investe, que consagra --- e a Igreja que acolhe e reconhece este dom, celebrando-o sacramentalmente. A vocação precede a ordenação; e a ordenação confere o que a vocação prometia, mas não podia por si só realizar.

\section{O sacerdócio na história da salvação (1539--1553)}

Para compreender o sacerdócio ministerial do Novo Testamento, o Catecismo recorre à pedagogia histórica de Deus: o sacerdócio da Antiga Aliança, com suas luzes e suas limitações, é a preparação e a prefiguração do sacerdócio único e definitivo de Cristo.

O povo escolhido foi constituído por Deus como ``um reino de sacerdotes e uma nação santa'' (Ex 19, 6). Mas, dentro de Israel, Deus separou uma das doze tribos, a de Levi, para o serviço litúrgico. Um rito especial consagrava os sacerdotes. Postos a agir ``em favor dos homens em relação a Deus, para oferecer dons e sacrifícios pelos pecados'' (Hb 5, 1), estes sacerdotes, porém, eram incapazes de alcançar uma santificação definitiva: seus sacrifícios eram repetidos incessantemente e não podiam realizar a salvação. Somente o sacrifício de Cristo a realizaria.

A liturgia da Igreja vê nesse sacerdócio antigo --- no de Aarão, no serviço dos levitas, na instituição dos setenta anciãos --- uma prefiguração do ministério ordenado do Novo Testamento. Por isso, nas orações consecratórias de bispos, presbíteros e diáconos, a Igreja retoma explicitamente esta memória: nomeia os sacerdotes do Antigo Testamento, o espírito comunicado aos setenta sábios, a plenitude do poder transmitida aos filhos de Aarão, os levitas chamados ao serviço do tabernáculo. Esta anamnese litúrgica não é mera arqueologia religiosa: é o reconhecimento de que a história da salvação tem uma coerência interior, uma única lógica divina que desemboca em Cristo.

Com efeito, tudo o que o sacerdócio da Antiga Aliança prefigurava encontra seu cumprimento em Cristo Jesus, ``único mediador entre Deus e os homens'' (1 Tm 2, 5). Cristo é o ``sumo sacerdote segundo a ordem de Melquisedeque'' (cf.~Hb 5, 6-10; 7, 11; Sl 110, 4): ``santo, inculpável, imaculado'' (Hb 7, 26), que ``por uma única oblação aperfeiçoou para sempre os que são santificados'' (Hb 10, 14). O sacrifício redentor de Cristo é único, realizado de uma vez para sempre; mas torna-se presente no sacrifício eucarístico da Igreja. O mesmo vale para o único sacerdócio de Cristo: ele se torna presente pelo sacerdócio ministerial, sem que a unicidade do sacerdócio de Cristo seja jamais diminuída. Como diz uma bela frase da tradição: ``Somente Cristo é o verdadeiro sacerdote; os outros são apenas seus ministros.''

Cristo fez da Igreja ``um reino de sacerdotes para o seu Deus e Pai'' (Ap 1, 6). Assim, toda a comunidade dos fiéis é, enquanto tal, sacerdotal. Os fiéis exercem seu sacerdócio batismal participando, cada qual conforme sua vocação, da missão de Cristo sacerdote, profeta e rei. Pelos sacramentos do Batismo e da Confirmação, os fiéis são ``consagrados a formar um sacerdócio santo'' (1 Pd 2, 5). Mas o sacerdócio ministerial ou hierárquico dos bispos e dos sacerdotes e o sacerdócio comum de todos os fiéis participam, ``cada um do seu modo próprio, do único sacerdócio de Cristo''. Ainda que ``ordenados uns aos outros'', eles diferem essencialmente: enquanto o sacerdócio comum se exerce pelo desenvolvimento da graça batismal --- uma vida de fé, de esperança, de caridade, uma vida segundo o Espírito ---, o sacerdócio ministerial está a serviço do sacerdócio comum e se destina ao desenvolvimento da graça batismal de todos os cristãos.

No serviço eclesial do ministro ordenado, é o próprio Cristo que está presente em sua Igreja como Cabeça de seu Corpo, Pastor do seu rebanho, sumo sacerdote do sacrifício redentor, Mestre da Verdade. É isso que a Igreja quer dizer quando afirma que o sacerdote, em virtude do sacramento da Ordem, age \textit{in persona Christi Capitis}: é o mesmo sacerdote, Cristo Jesus, cuja pessoa sagrada seu ministro representa verdadeiramente. Contudo --- e o Catecismo é preciso quanto a este ponto --- esta presença de Cristo no ministro não significa que o ministro esteja preservado de toda fragilidade humana, do espírito de dominação, do erro ou mesmo do pecado. O poder do Espírito Santo não garante da mesma maneira todos os atos dos ministros: ora garante os sacramentos, de modo que mesmo o pecado do ministro não pode impedir o fruto da graça; mas em muitos outros atos o ministro deixa traços humanos que nem sempre são sinais de fidelidade ao Evangelho.

O sacerdócio ministerial é serviço (\textit{ministerium}): ``o encargo que o Senhor confiou aos pastores do seu povo é, no sentido estrito do termo, um serviço'' (LG 24). Está inteiramente ordenado a Cristo e aos homens. Depende totalmente de Cristo e do seu sacerdócio único, e foi instituído para o bem dos homens e a comunhão da Igreja. O ministério ordenado não representa a si mesmo, mas representa Cristo diante da assembleia; e, ao mesmo tempo, representa toda a Igreja diante de Deus, quando apresenta ao Pai a oração da Igreja e, sobretudo, quando oferece o sacrifício eucarístico. ``Em nome de toda a Igreja'' não significa que os sacerdotes sejam delegados da comunidade: a oração e a oblação da Igreja são inseparáveis da oração e da oblação de Cristo, sua Cabeça; e é sempre Ele que adora em e pela Igreja. É precisamente porque o sacerdócio ministerial representa Cristo que ele pode representar a Igreja.

\section{Os três graus do sacramento da Ordem (1554--1571)}

O ministério eclesiástico de instituição divina é exercido em diferentes graus por aqueles que, desde os tempos mais antigos, são chamados de bispos, sacerdotes e diáconos. A doutrina católica, expressa na liturgia, no Magistério e na prática constante da Igreja, reconhece que há dois graus de participação ministerial no sacerdócio de Cristo --- o episcopado e o presbiterado --- e um grau de serviço, o diaconato, destinado a assistir e servir os primeiros. Por isso, o termo \textit{sacerdos}, no uso corrente, designa bispos e sacerdotes, mas não os diáconos. Contudo, a doutrina católica ensina que os três graus são conferidos por um ato sacramental chamado ``ordenação'', isto é, pelo sacramento da Ordem.

\subsection*{O episcopado}

Entre as diversas funções exercidas na Igreja desde os primeiros tempos, o lugar principal cabe, segundo o testemunho da Tradição, àqueles que, pela sua nomeação à dignidade e à responsabilidade de bispo, e em virtude da sucessão ininterrupta que remonta ao início, são considerados como transmissores da linha apostólica. Para cumprir sua missão exaltada, ``os apóstolos foram dotados por Cristo com uma efusão especial do Espírito Santo e, pela imposição das mãos, transmitiram a seus auxiliares o dom do Espírito, que é transmitido até nosso dia pela consagração episcopal'' (LG 21).

O Concílio Vaticano~II ensina que ``a plenitude do sacramento da Ordem é conferida pela consagração episcopal''. Esta plenitude, que a tradição litúrgica e os Pais da Igreja chamam de sumo sacerdócio e ápice do ministério sagrado, confere juntamente com o encargo de santificar também os encargos de ensinar e de governar. Pelo dom do Espírito Santo e a impressão de um caráter sagrado, os bispos tomam, de maneira eminente e visível, o lugar do próprio Cristo --- professor, pastor e sacerdote --- e agem como seus representantes (\textit{in Eius persona agant}).

A natureza colegial da ordem episcopal manifesta-se, entre outras coisas, na prática antiga da Igreja de chamar vários bispos para participar da consagração de um novo bispo. Nos nossos dias, a ordenação legítima de um bispo requer a intervenção especial do Bispo de Roma, que é o vínculo visível supremo da comunhão das Igrejas particulares na única Igreja e o garante de sua liberdade. Na condição de vigário de Cristo, cada bispo tem a solicitude pastoral da Igreja particular que lhe foi confiada; mas ao mesmo tempo, colegialmente com todos os seus irmãos no episcopado, partilha a solicitude por todas as Igrejas: cada bispo legítimo, como legítimo sucessor dos apóstolos, é, por instituição e preceito divino, corresponsável com os demais bispos pela missão apostólica da Igreja.

\subsection*{O presbiterado}

Cristo, que o Pai santificou e enviou ao mundo, comunicou aos bispos, por meio dos Apóstolos, a participação em sua consagração e missão; e estes, por sua vez, confiaram legitimamente, em graus variados, os membros da Igreja o exercício do ministério. ``A função do ministério dos bispos foi transmitida, em grau subordinado, aos sacerdotes, para que, constituídos na ordem do sacerdócio, fossem cooperadores da ordem episcopal para o desempenho reto da missão apostólica confiada por Cristo'' (PO 2). Os sacerdotes participam assim das dimensões universais da missão que Cristo confiou aos Apóstolos: o dom espiritual que recebem na ordenação não os prepara para uma missão limitada, mas para ``a missão de salvação mais plena, até os confins da terra'' (PO 10).

O sacerdócio dos sacerdotes --- pressupondo os sacramentos da iniciação --- é conferido pelo seu próprio sacramento particular. Pela unção do Espírito Santo, os sacerdotes são marcados com um caráter especial e assim configurados a Cristo Sacerdote, de tal modo que podem agir na pessoa de Cristo Cabeça. É na celebração eucarística que exercem de modo supremo o seu ministério sagrado: aí, agindo na pessoa de Cristo e proclamando o seu mistério, unem as oblações dos fiéis ao sacrifício de seu Cabeça, e no sacrifício da Missa tornam presente e atualizam o único sacrifício do Novo Testamento, o de Cristo que se ofereceu a si mesmo como vítima imaculada ao Pai. Deste único sacrifício, todo o ministério sacerdotal retira a sua força.

Os sacerdotes são cooperadores prudentes da ordem episcopal e seu suporte; juntamente com seu bispo formam um único corpo sacerdotal (\textit{presbyterium}) dedicado a diversas tarefas. Em cada assembleia local de fiéis, representam em certo sentido o bispo, com quem se associam com toda a confiança e generosidade, assumindo parte de suas atribuições e cuidados. A promessa de obediência ao bispo feita no momento da ordenação e o ósculo de paz recebido ao final da liturgia da ordenação expressam que o bispo os acolhe como seus cooperadores, seus filhos, seus irmãos e seus amigos, e que eles lhe devem amor e obediência. ``Todos os sacerdotes constituídos na ordem do sacerdócio pelo sacramento da Ordem estão unidos entre si por uma íntima fraternidade sacramental'' (PO 8), encontrando na unidade do \textit{presbyterium} a expressão litúrgica mais concreta da sua comunhão mútua.

\subsection*{O diaconato}

Os diáconos recebem a imposição das mãos ``não para o sacerdócio, mas para o serviço''. No Rito Romano, são ordenados pelo bispo sozinho. O texto da prece consecratória no rito de ordenação dos diáconos evoca a instituição de Estêvão e dos ``sete homens cheios do Espírito e de sabedoria'' (cf.~At 6, 3-6), que foram estabelecidos ao serviço da comunidade, e também os levitas do Antigo Testamento escolhidos para o serviço do tabernáculo. ``Fortalecidos pela graça sacramental, os diáconos servem o Povo de Deus, em comunhão com o bispo e seu presbitério, no serviço (\textit{diakonia}) da liturgia, da palavra e da caridade'' (LG 29). O Concílio Vaticano~II restabeleceu o diaconato permanente como um grau próprio e estável da hierarquia, que pode ser conferido também a homens casados.

\section{A celebração do sacramento e seu ministro (1572--1580)}

Dada a importância que a ordenação de um bispo, de um sacerdote ou de um diácono tem para a vida da Igreja particular, sua celebração reclama a participação do maior número possível de fiéis. Deve ocorrer preferencialmente no domingo, na catedral, com a solenidade adequada à circunstância. As três ordenações --- do bispo, do sacerdote e do diácono --- seguem o mesmo movimento e têm o seu lugar próprio dentro da liturgia eucarística.

O rito essencial do sacramento da Ordem, para os três graus, consiste na imposição das mãos do bispo sobre a cabeça do ordenando e na oração consecratória específica do bispo, pedindo a Deus a efusão do Espírito Santo e de seus dons próprios ao ministério para o qual o candidato está sendo ordenado. Como em todos os sacramentos, ritos adicionais cercam a celebração, variando bastante entre as diferentes tradições litúrgicas, mas tendo em comum o objetivo de exprimir os múltiplos aspectos da graça sacramental. No Rito Latino, os ritos iniciais --- apresentação e eleição do ordenando, instrução pelo bispo, exame do candidato, ladainha dos santos --- atestam que a escolha do candidato foi feita segundo a prática da Igreja e preparam para o ato solene da consagração. Seguem-se ritos que exprimem e completam simbolicamente o mistério realizado: a unção com o santo crisma para bispos e sacerdotes (sinal da unção especial do Espírito Santo que fecunda o ministério); a entrega do Evangelho, do anel, da mitra e do báculo ao bispo; a apresentação ao sacerdote da patena e do cálice com ``a oferta do santo povo''; a entrega do Evangelho ao diácono encarregado de proclamar o Evangelho de Cristo.

O ministro do sacramento da Ordem é o bispo. ``Uma vez que o sacramento da Ordem é o sacramento do ministério apostólico, compete aos bispos, como sucessores dos apóstolos, transmitir o ``dom do Espírito''... Bispos validamente ordenados --- isto é, aqueles que estão na linha da sucessão apostólica --- conferem validamente os três graus do sacramento da Ordem'' (cf.~DS 794; 802; CIC, can. 1012). Cristo mesmo escolheu os Apóstolos e lhes deu participação em sua missão e autoridade. Elevado à direita do Pai, não abandonou o seu rebanho, mas continua a conduzi-lo por meio dos bispos que prosseguem hoje a sua obra. É Cristo que age pelos bispos ordenados.

A disciplina sobre o candidato à Ordem apresenta dois pontos que o Catecismo expõe com clareza. Em primeiro lugar, ``somente um homem batizado (\textit{vir}) recebe validamente a sagrada ordenação'' (CIC, can. 1024): o Senhor Jesus escolheu homens para formar o colégio dos Doze, e os Apóstolos fizeram o mesmo ao escolher colaboradores para os suceder em seu ministério. O colégio dos bispos, com quem os sacerdotes estão unidos no sacerdócio, faz do colégio dos Doze uma realidade sempre presente e sempre ativa até a volta de Cristo. A Igreja reconhece-se vinculada por esta escolha feita pelo próprio Senhor; razão pela qual não é possível a ordenação de mulheres. Em segundo lugar, ninguém tem direito de receber o sacramento da Ordem: ninguém o reivindica para si; é chamado por Deus. Quem julga reconhecer os sinais de uma chamada de Deus ao ministério ordenado deve submeter com humildade seu desejo à autoridade da Igreja, que tem a responsabilidade e o direito de chamar alguém a receber as ordens. Como toda graça, este sacramento só pode ser recebido como dom imerecido.

Quanto ao celibato, todos os ministros ordenados da Igreja Latina --- com exceção dos diáconos permanentes --- são normalmente escolhidos entre homens de fé que vivem em celibato e que têm intenção de permanecer célibes ``por causa do Reino dos Céus'' (Mt 19, 12). Chamados a consagrar-se com coração indiviso ao Senhor e ``às coisas do Senhor'' (1 Cor 7, 32), eles se dão inteiramente a Deus e aos homens. O celibato é um sinal desta vida nova ao serviço da qual o ministro da Igreja é consagrado; acolhido com coração alegre, proclama radiantemente o Reino de Deus. Nas Igrejas do Oriente, vigora uma disciplina diferente: embora os bispos sejam escolhidos unicamente entre os celibatários, homens casados podem ser ordenados diáconos e sacerdotes --- prática há muito reconhecida como legítima e que naquelas comunidades produz frutos de fecundidade pastoral. O celibato sacerdotal, no entanto, é muito estimado também nas Igrejas orientais, e muitos sacerdotes escolheram-no livremente por causa do Reino de Deus. Tanto no Oriente quanto no Ocidente, quem já recebeu o sacramento da Ordem não pode mais contrair matrimônio.

\section{Os efeitos do sacramento da Ordem (1581--1589)}

O sacramento da Ordem configura o ordenando a Cristo por uma graça especial do Espírito Santo, para que sirva de instrumento de Cristo para a sua Igreja. Pela ordenação, é habilitado a agir como representante de Cristo Cabeça da Igreja, no seu tríplice ofício de sacerdote, profeta e rei.

Como o Batismo e a Confirmação, esta participação no ofício de Cristo é concedida de uma vez por todas. O sacramento da Ordem, como os outros dois, imprime um caráter espiritual indelével e não pode ser repetido nem conferido temporariamente. É verdade que alguém validamente ordenado pode, por justa causa, ser dispensado das obrigações e funções ligadas à ordenação, ou ter-lhes o exercício proibido; mas não pode voltar a ser leigo no sentido estrito, porque o caráter impresso pela ordenação é para sempre. A vocação e missão recebidas no dia da sua ordenação marcam-no permanentemente.

Mas --- e o Catecismo faz aqui uma afirmação de grande alcance pastoral e teológico --- como é, em última análise, Cristo quem age e realiza a salvação pelo ministro ordenado, a indignidade eventual do ministro não impede Cristo de agir. Santo Agostinho formula esta verdade com vigor: o ministro orgulhoso deve ser ranqueado com o diabo, mas o dom de Cristo não é por isso profanado: ``o que flui por meio dele conserva a sua pureza, e o que passa por ele permanece intacto e chega à terra fértil''. O poder espiritual do sacramento é comparável à luz: os que são iluminados por ele recebem-no em toda a sua pureza; e se ele passou por seres manchados, não se manca. Esta objetividade da graça sacramental é fundamento de segurança para os fiéis e de humildade para os ministros.

A graça especial do sacramento da Ordem é a configuração com Cristo sacerdote, mestre e pastor, de quem o ordenado é feito ministro. Para o bispo, é antes de tudo uma graça de força --- ``o espírito de governo'', segundo a prece de consagração episcopal do Rito Latino ---: a graça de guiar e defender a sua Igreja com força e prudência de pai e pastor, com amor gratuito por todos e amor preferencial pelos pobres, pelos doentes e pelos necessitados. Esta graça leva-o a proclamar o Evangelho a todos, a ser o modelo do seu rebanho, a caminhar à sua frente no caminho da santificação, a identificar-se na Eucaristia com Cristo sacerdote e vítima, sem temer dar a vida pelas suas ovelhas. Para os presbíteros, o dom espiritual conferido pela ordenação os habilita a agir na pessoa de Cristo Cabeça: proclamar o Evangelho, apascentar os fiéis, celebrar o culto divino, oferecer o sacrifício eucarístico em nome de todo o povo. Para os diáconos, ``fortalecidos pela graça sacramental, são dedicados ao Povo de Deus, em conjunto com o bispo e seu presbitério, no serviço da liturgia, do Evangelho e das obras de caridade'' (LG 29).

Diante da grandeza da graça e do ministério sacerdotais, os santos doutores sentiram um urgente apelo à conversão, para conformar toda a sua vida àquele de quem o sacramento os fez ministros. São Gregório Nazianzeno, ainda jovem sacerdote, exclamava: ``Precisamos começar por nos purificar antes de purificar os outros; precisamos ser instruídos para poder instruir, tornarmo-nos luz para iluminar, aproximarmo-nos de Deus para levá-L'O aos outros, ser santificados para santificar, conduzir pela mão e aconselhar prudentemente...'' E o santo Cura d'Ars: ``O sacerdote continua a obra da Redenção na terra... O sacerdócio é o amor do Coração de Jesus.'' Tais testemunhos não são hagiografia piegas: são descrições do que a graça da Ordem opera em quem a acolhe com docilidade, deixando Cristo ser, de fato, o único Sacerdote que age por meio do ministro.

\section{Em síntese (1590--1600)}

São Paulo dizia ao seu discípulo Timóteo: ``Recordo-te que deves reavivar o dom de Deus que está em ti pela imposição das minhas mãos'' (2 Tm 1, 6) e ``Se alguém aspira ao episcopado, deseja uma obra excelente'' (1 Tm 3, 1). A Tito escrevia: ``Por isso te deixei em Creta, para que regulasses o que ainda faltava e constituísses presbíteros em cada cidade, conforme te ordenei'' (Tt 1, 5). Com estas palavras apostólicas, o Catecismo abre sua síntese doutrinária sobre o sacramento da Ordem, recordando que não é uma invenção eclesiástica tardia, mas um dado original da fé, enraizado na Escritura e vivido pela Igreja desde o seu início.

A Igreja toda é um povo sacerdotal. Pelo Batismo, todos os fiéis participam do sacerdócio de Cristo. Esta participação é chamada ``sacerdócio comum dos fiéis''. Fundado neste sacerdócio comum e ordenado ao seu serviço, existe outra participação na missão de Cristo: o ministério conferido pelo sacramento da Ordem, cujo encargo é servir em nome e na pessoa de Cristo Cabeça no meio da comunidade. O ministério ordenado é essencialmente diferente do sacerdócio comum, não em grau, mas em natureza: o sacerdote não é ``mais cristão'' que os leigos; é um tipo diferente de presença de Cristo na Igreja.

O sacerdócio ministerial é exercido em três graus: o episcopado, o presbiterado e o diaconato. Os três são conferidos pelo sacramento da Ordem, dado pela imposição das mãos e pela oração consecratória. Os bispos são os sucessores dos Apóstolos; recebem a plenitude do sacramento da Ordem; são os guardiães autênticos da fé, os pastores das Igrejas particulares e os membros do colégio episcopal, que tem como cabeça o Bispo de Roma. Os presbíteros são os cooperadores dos bispos; compartilham do seu sacerdócio e da sua missão, formando com eles o \textit{presbyterium} da Igreja particular. Os diáconos são ordenados não para o sacerdócio, mas para o serviço: assistem os bispos e os sacerdotes na celebração dos sacramentos, na proclamação do Evangelho e nas obras de caridade.

Somente o bispo é ministro do sacramento da Ordem. Somente um homem batizado pode recebê-lo validamente. Ninguém tem direito a ele: é uma vocação e um dom. O sacramento da Ordem imprime um caráter indelével: quem o recebe é configurado a Cristo para sempre, mesmo que seja dispensado do exercício das funções que lhe estão ligadas. A dignidade pessoal do ministro não é condição para a validade dos sacramentos: a graça de Cristo permanece pura mesmo quando passa por um ministro indigno.

A graça própria do sacramento da Ordem é a configuração com Cristo sacerdote, profeta e pastor, e a missão de servir o Povo de Deus em nome de Cristo Cabeça. Para o bispo, esta graça é força de governo, prudência pastoral, coragem apostólica e amor pelos pobres. Para o sacerdote, é a habilitação para agir \textit{in persona Christi Capitis}, especialmente na celebração da Eucaristia. Para o diácono, é o fortalecimento para o serviço da liturgia, da Palavra e da caridade. Em todos os três graus, o sacramento da Ordem perpetua a missão que Cristo confiou a seus Apóstolos e que os Apóstolos transmitiram a seus sucessores: anunciar, celebrar e servir --- até que Ele venha.
